\id{ҒТАМР 65.59.31}{}

\begin{header}
\swa{Ш.Б.~Байтукенова, M.~Бақытбек, С.Б.~Байтукенова, А.Т.~Костанова}{ПІСІРІЛГЕН ШҰЖЫҚ ӨНДІРІСІНДЕ АҚУЫЗ-МАЙ ЭМУЛЬСИЯСЫНЫҢ ТЕХНОЛОГИЯСЫН ЖАСАУ}

\tsp{1}Ш.Б.~Байтукенова\envelope,
\tsp{1}M.~Бақытбек,
\tsp{2}С.Б.~Байтукенова,
\tsp{1}А.Т.~Костанова
\end{header}

\begin{affil}
\tsp{1}«С.Сейфуллин атындағы Қазақ агротехникалық зерттеу университеті» КеАҚ, Астана, Қазақстан,

\tsp{2}«Қ.Құлажанов атындағы Қазақ технология және бизнес университеті» АҚ, Астана, Қазақстан

\corrauthor{Корреспондент-автор: baytukenova75@mail.ru}
\end{affil}

Бұл мақалада ақуыз-май эмульсиясының технологиясын жасау және оны
пісірілген шұжық өнімдерінде қолдану тиімділігі зерттелді. Зерттеудің
негізгі мақсаты -- ақуыз-май эмульсияның функционалдық, құрылымдық және
органолептикалық қасиеттеріне әсерін анықтау және тағамдық өнімдердің
сапасын арттыруға бағытталған технологиялық шешімдерді жасау болып
табылады.

Зерттеу барысында бірнеше тәжірибелік үлгілер дайындалып, олардың
құрылымдық-механика\-лық және органолептикалық көрсеткіштері және
энергетикалық құндылығы қарастырылды. Алынған нәтижелер компоненттердің
үйлесімділігі эмульсияның тұрақтылығы мен біркелкі құрылымына, сондай-ақ
дайын өнімнің дәмдік қасиеттерінің жақсаруына тікелей әсер ететінін
көрсетті. Сонымен қатар, ақуыз-май эмульсиясының құрамындағы әр
компоненттің өзара қатынасы өнімнің технологиялық өңдеуге төзімділігін
қамтамасыз етіп, тағамдық құндылығын сақтауға мүмкіндік береді.

Пісірілген шұжық рецептурасы құрамына ақуыз-май эмульсиясын 15\% енгізу
олардың құрылымын және дәмдік қасиетін жақсартады, сондай-ақ текстурасын
жұмсартатыны анықталды. Эмульсия өнімнің ылғал ұстағыш қасиетін
арттырып, 4 °C температурада сақтаған кезде бақылау үлгісімен
салыстырғанда 5-7 күнге ұзарды, бұл әсіресе дайын өнімнің сапасын сақтау
үшін маңызды. Жүргізілген зерттеулер нәтижесінде пісірілген шұжық
рецептурасы құрамына ақуыз-май эмульсиясын 15\% енгізу дайын өнімнің
құрылымдық, функционалдық және органолептикалық қасиеттерін жақсартуға
мүмкіндік беріп қана қоймай, өндірістік тиімділікті арттырады және
заманауи тұтынушылардың тағам өнімдеріне қойылатын талаптарына толық
сәйкес келеді.

{\bfseries Түйін сөздер:} соя сүті, зәйтүн майы, соя изоляты, ақуыз-май
эмульсиясы, физика-химиялық қасиеттері, құрылымдық-механикалық
қасиеттері, пісірілген шұжық өнімдері, технологиялық сұлба.

\begin{header}
РАЗРАБОТКА ТЕХНОЛОГИИ БЕЛКОВО-ЖИРОВОЙ ЭМУЛЬСИИ В ПРОИЗВОДСТВЕ ВАРЕНОЙ КОЛБАСЫ

\tsp{1}Ш.Б.~Байтукенова\envelope,
\tsp{1}М.~Бақытбек,
\tsp{2}С.Б.~Байтукенова,
\tsp{1}А.Т.~Костанова
\end{header}

\begin{affil}
\tsp{1}НАО «Казахский агротехнический, исследовательский университет им. С.Сейфуллина», Астана, Казахстан,

\tsp{2}АО «Казахский университет технологии и бизнеса, им. К. Кулажанова» Астана, Казахстан,

e-mail: baytukenova75@mail.ru
\end{affil}

В данной статье изучена эффективность разработки технологии
белково-жировой эмульсии \\(БЖЭ) и ее применения в производстве вареной
колбасы. Основной целью исследования является выявление влияния
белково-жировой эмульсии на функциональные, структурные и
органолептические свойства и разработка технологических решений,
направленных на повышение качества пищевых продуктов.

В ходе исследования было подготовлено несколько опытных образцов БЖЭ,
рассмотрены их структурно-механические и органолептические показатели и
энергетическая ценность. Полученные результаты показали, что
совместимость компонентов напрямую влияет на стабильность и однородную
структуру эмульсии, а также на улучшение вкусовых свойств готового
продукта. Кроме того, соотношение каждого компонента в составе
белково-жировой эмульсии позволяет сохранить пищевую ценность,
обеспечивая устойчивость продукта к технологической обработке.

Установлено, что введение 15\% белково-жировой эмульсии в рецептуру
вареной колбасы улучшает их структуру и вкусовые качества, а также
смягчает текстуру. Эмульсия увеличивает влагоудерживающие свойства
продукта и позволяет продлить срок хранения на 5-7 дней по сравнению с
контрольным образцом при температуре 4 °C, что особенно важно для
поддержания качества готового продукта. В результате проведенных
исследований введение 15\% белково-жировой эмульсии в рецептуру вареной
колбасы не только позволило улучшить структурные, функциональные и
органолептические свойства готового продукта, но и повысило
эффективность производства и полностью соответствовало требованиям
современных потребителей к пищевой продукци

{\bfseries Ключевые слова:} соевое молоко, оливковое масло, соевый изолят,
белково-жировая эмульсия, физико-химические свойства,
структурно-механические свойства, вареные колбасные изделия,
технологическая схема.

\begin{header}
DEVELOPMENT OF PROTEIN-FAT EMULSION TECHNOLOGY IN THE PRODUCTION OF BOILED SAUSAGE

\tsp{1}Sh.B.~Baitukenova\envelope,
\tsp{1}M.~Bakytbek,
\tsp{2}S.B.~Baitukenova,
\tsp{1}А.Т.~Kostanova
\end{header}

\begin{affil}
\tsp{1}S. Seifullin Kazakh Agrotechnical Research, University, Astana, Kazakhstan,

\tsp{2}«K. Kulazhanov Kazakh University of Technology, and Business, Astana, Kazakhstan,

e-mail: baytukenova75@mail.ru
\end{affil}

This article examines the effectiveness of the development of
protein-fat emulsion (BGE) technology and its application in the
production of boiled sausage. The main purpose of the study is to
identify the effect of protein-fat emulsion on functional, structural
and organoleptic properties and to develop technological solutions aimed
at improving food quality.

In the course of the study, several prototypes of BZHE were prepared,
their structural, mechanical and organoleptic characteristics and energy
value were considered. The results showed that the compatibility of the
components directly affects the stability and homogeneous structure of
the emulsion, as well as the improvement of the taste properties of the
finished product. In addition, the ratio of each component in the
protein-fat emulsion makes it possible to preserve the nutritional
value, ensuring the product' s resistance to
technological processing.

It was found that the introduction of 15\% protein-fat emulsion into the
recipe of boiled sausage improves their structure and taste, as well as
softens the texture. The emulsion increases the moisture-retaining
properties of the product and allows you to extend the shelf life by 5-7
days compared to the control sample at a temperature of 4 °C, which is
especially important for maintaining the quality of the finished
product. As a result of the research, the introduction of 15\%
protein-fat emulsion into the recipe of boiled sausage not only improved
the structural, functional and organoleptic properties of the finished
product, but also increased production efficiency and fully met the
requirements of modern consumers for food products.

{\bfseries Keywords:} soy milk, olive oil, soy isolate, protein-fat
emulsion, physico-chemical properties, structural and mechanical
properties, boiled sausages, technological scheme.

\begin{multicols}{2}
{\bfseries Кіріспе.} Соңғы жылдары Қазақстанның тамақ өнеркәсібінде өсімдік
текті ингредиенттерді қолдану қарқынды дамып келеді. Тұтынушылардың
табиғи, қауіпсіз және құрамы жағынан теңгерімді өнімдерге деген сұранысы
артқан сайын, өңделген өсімдік шикізатына негізделген технологиялардың
маңызы барған сайын артып отыр. Әсіресе соя дәнінен алынатын ақуыздар
мен майлар тағам өндірісінде кеңінен қолданылатын функционалды компонент
ретінде ерекше қызығушылық тудыруда.

Қазақстан Республикасы Стратегиялық жоспарлау және реформалар
агенттігінің 2025 жылғы мәліметтері бойынша, елде майлы дақылдар
өндірісі соңғы бес жылда 35-40\%-ға артқан, оның ішінде соя дақылының
үлесі тұрақты өсім көрсетуде.2025 жылы Республикада 190 мың тоннадан
астам соя жиналып, бұл оған өңдеу саласында жаңа технологияларды
енгізуге жағдай жасады. Сонымен қатар, соңғы үш жылда өсімдік
негізіндегі өнімдерге сұраныс 25\%-ға артқаны байқалады, бұл халықтың
дұрыс тамақтану мәдениетінің қалыптасуына және функционалды өнімдерге
бет бұруына тікелей байланысты {[}1{]}.

Соя дәнінің құрамындағы 40\%-ға дейінгі ақуыз, 18-20\% май, маңызды
аминқышқылдары, фосфолипидтер және биологиялық белсенді қосылыстар сояны
тағамдық эмульсиялар алу үшін тиімді шикізатқа айналдырады. Соя
ақуызының табиғи эмульгирлеуші қабілеті, май фазасын тұрақтандыру
қасиеті және гидратациялық мүмкіндігі ақуыз -- май қоспаларын дайындауда
үлкен артықшылық береді. Мұндай эмульсиялар құрылым түзуші, тұтқырлық
қалыптастырушы, майды біркелкі таратуға көмектесетін технологиялық
компонент ретінде жоғары бағаланады {[}2{]}.

Соңғы зерттеулер соя ақуызы негізіндегі эмульсиялардың тағам құрылымының
тұрақтылығын арттырып, ылғал және май ұстағыштық мүмкіндігін күшейтіп,
өнімнің органолептикалық қасиеттерін жақсартатынын көрсетіп отыр.
Эмульсия құрамындағы ақуыз-май қатынасының дұрыс таңдалуы дайын өнімнің
консистенциясына, тұтқырлығына, энергетикалық құндылығына және сақтау
тұрақтылығына тікелей әсер етеді {[}3{]}.

Қазақстанда өсімдік ингредиенттері өндірісінің артуы - соя негізіндегі
ақуыз-май эмульсияларын кеңінен енгізу үшін қолайлы жағдай
қалыптастырады. Бұл технологияны қолдану: өнімнің құрылымдық және
технологиялық қасиеттерін жақсартады; табиғи эмульгирлеуші компоненттер
есебінен қоспалардың мөлшерін азайтады; функционалды құндылығын
арттырады; өндіріске жұмсалатын шығынды төмендетеді; отандық
кәсіпорындардың бәсекеге қабілеттілігін күшейтеді.

{\bfseries Материалдар мен әдістер.} Зерттеу объектісі:~ақуыз-май қоспасы,
соя сүті, соя ақуызының изоляты, өсімдік майы - зәйтүн майы.

Шикізаттың сапалық көрсеткіштері: соя сүті - Қазақстандағы «АгроСоя» ЖШС
өндірген, 3-4 \% ақуыз үлесімен, ГОСТ 32252-2013 стандартына сәйкес; соя
ақуызының изоляты -- Шаньсун-90 маркасы (Shansong 90), 90\% ақуыз
үлесімен; зәйтүн майы - Италиядан импортталған «Bertolli Original»
бірінші суық сығымдалған, экстра-класс, қаныққан май қышқылдары 14-15\%,
моноқанықпаған 70-78\%. Физика-химиялық сипаттамалары: соя сүті -- рН
6,5-7,0, ылғал 90\%; соя изоляты - гидратация қабілеті 1:5; зәйтүн
майының тотығу индексі <0,5 мэкв/кг ((миллививаленттерінде).

Зерттеулер С.Сейфуллин атындағы Қазақ агротехникалық зерттеу
университетінің етті қайта өңдеу эксперименттік-өндірістік цехта және
зертхана жағдайында жүргізілді. Өнімдегі ылғалдың, ақуыздың, майдың,
минералды заттардың массалық үлесі стандартты әдістер бойынша анықталды.
Органолептикалық көрсеткіштерін анықтауда сынамалар 5 балдық жүйе
бойынша сыртқы түрі, иісі, дәмі, консистенциясы және түсі бойынша
бағаланды. Органолептикалық бағалау ГОСТ 9959-2015 сәйкес жүргізілді.

Су байланыстырғыш және май сіңіру қабілеттілігі (СБҚ және МСҚ), рН
көрсеткіштері мен шығымы анықталды (нақты нәтижелер төмендегі кестеде).
Микробиологиялық көрсеткіштер ГОСТ 54354-2011 сәйкес зерттелді: жалпы
бактерия саны <103 КОЕ/г, колиформалар жоқ, сақтау мерзімі
4°C-та 14-21 күн (бақылау үлгісі 7-14 күн). Зерттеу нәтижелері
дисперсиялық талдау (ANOVA) әдісімен статистикалық өңдеуден өткізілді.

{\bfseries Нәтижелер мен талқылау.} Ақуыз-май эмульсиясының өндірістік және
функционалды қасиеттерін зерттеу қазіргі тағам өнеркәсібінде маңызды
мәселе болып табылады. Бұл қоспаның негізгі мақсаты ақуыз бен майдың
біркелкі таралуын қамтамасыз ету, тұрақты құрылымға ие эмульсия алу,
ылғал мен майды ұстап тұру қабілетін жақсарту, сондай ақ тағамдық
қасиеттерін сақтау {[}4{]}.

Қоспаның сапасы бірнеше факторға байланысты: шикізат құрамындағы ақуыз
бен майдың қатынасы, эмульсияның құрылымдық тұрақтылығы, су
байланыстырғыш қабілеті және органолептикалық сипаттар (түс,
консистенция, дәм, және иіс). Осы факторлардың оңтайлы үйлесімі
ақуыз-май эмульсиясының технологиялық және тағамдық қасиеттерін
максималды түрде арттыруға мүмкіндік береді {[}5{]}.

Ақуыз-май қоспасын алу үшін бірнеше өсімдік шикізаттары қарастырылды
соя, жасымық, ноқат және ұрмебұршақ. Әрқайсының өзіндік ерекшеліктері
бар және олардың химиялық құрамы мен функционалды қасиеттері қоспаның
сапасына тікелей әсер етеді. Соя ең жоғарғы ақуыз мөлшерімен
ерекшеленеді, сонымен қатар эмульсия түзу қабілеті жоғары, бұл оны
ақуыз-май эмульсиясының негізгі компоненті ретінде пайдалануға мүмкіндік
береді. Жасымық пен ноқат құрамында ақуыз жеткілікті және олар қоспаның
органолептикалық қасиеттерін жақсартуға, дәмін байытуға, құрылымын
жұмсартып, консистенциясы біркелкі етуге ықпал етеді. Ұрмебұршақ
қосылған қоспалар фарштың ылғал ұстай алу қабілетін арттырып, құрылымдық
тұрақтылығын нығайтады {[}6{]}.

Соя, жасымық, ноқат және үрмебұршақ -- тағам өнеркәсібінде маңызды орын
алатын өсімдік текті шикізаттар, олар ақуызға, көмірсуға, биологиялық
белсенді қосылыстарға бай болуымен, сондай-ақ
функционалдық-технологиялық қасиеттерімен ерекшеленеді. Бұл дақылдар
ақуыз-май қоспасын байыту, құрылымын жақсартуға мүмкіндік береді
{[}7{]}.1-ші кестеде өсімдік шикізаттарының химиялық құрамы
көрсетілген.
\end{multicols}
\vspace{-0.6cm}
\tcap{1-кесте. Өсімдік шикізаттарының химиялық құрамы}
\begin{longtblr}[
  label = none,
  entry = none,
]{
  cells = {c},
  cells = {font = \small},
  row{1} = {font=\bfseries\small},
  hlines,
  vlines,
}
Көрсеткіштер                           & Соя & Жасымық & Ноқат & Ұрмебұршақ \\
Ақуыздың массалық үлесі, \%            & 40  & 26      & 19    & 21         \\
Майдың массалық үлесі, \%              & 18  & 1       & 6     & 2          \\
Көмірсулардың массалық үлесі, \%       & 30  & 48      & 50    & 55         \\
Минералды заттардың массалық үлесі, \% & 4.5 & 3       & 2.6   & 3          \\
Фосфолипидтердың массалық үлесі, \%    & 12  & 10      & 10    & 10         
\end{longtblr}
\vspace{-0.6cm}
\tcap{2-кесте. Өсімдік май шикізаттарының физика-химиялық құрамы}
\begin{longtblr}[
  label = none,
  entry = none,
]{
  cells = {c},
  cells = {font = \small},
  row{1} = {font=\bfseries\small},
  hlines,
  vlines,
}
Көрсеткіштері                     & Зәйтүн  & Күнбағыс & Зығыр \\
Қаныққан май қышқылдары, \%       & 14-15   & 10-11    & 9-10  \\
Витамин Е, мг/100                 & 14-20   & 40-50    & 20-30 \\
Моноқанықпаған май қышқылдары, \% & 70-78   & 19-25    & 18-22 \\
Полиқанықпаған май қышқылдары, \% & 12-34   & 65-70    & 65-72 \\
Омега-6 (линол қышқылы), \%       & 11-12   & 60–65    & 15-18 \\
Омега-3 (α-линолен қышқылы), \%   & 0.5-1   & 1        & 50–60 \\
Полифенолдар, мг/кг               & 150-300 & 20–50    & 50–80 \\
Стеролдар мг/100                  & 100-150 & 100-120  & 60-80 
\end{longtblr}
\vspace{-0.6cm}
\begin{multicols}{2}
Сояның басты артықшылығы жоғары ақуыз мөлшері мен функционалды
қасиеттері, бұл эмульсияның құрылымын тұрақтандырады, май мен судың
біркелкі бөлінуін қамтамассыз етеді және өнімнің органолептикалық
көрсеткіштерін жақсартады. Сонымен қатар, соя ақуызы биологиялық
құндылығы жоғары, құрамындағы аминқышқылдар және май қышқылдары
эмульсияның тағамдық қасиеттерін арттыруға септігін тигізеді. Сондықтан
ақуыз-май қоспасын алу үшін негізгі шикізат ретінде соя таңдалды {[}8,
9{]}.

Ақуыз-май эмульсиясының функционалды қасиеттерін арттыру және өнімнің
құрылымдық тұрақтылығын қамтамассыз ету үшін әртүрлі өсімдік майларыда
пайдаланылады. Зәйтүн, күнбағыс және зығыр майларының негізгі химиялық
құрамы 2-ші кестеде көрсетілген.

Зәйтүн, күнбағыс және зығыр майлары тағам өнеркәсібінде кеңінен
қолданылатын өсімдік майлары болып табылады және олардың технологиялық
тиімділігі тікелей химиялық құрамына байланысты анықталады. Бұл
майлардың негізгі компоненті -- триглицеридтер, бірақ май қышқылдарының
арақатынасы әрқайсысында айтарлықтай өзгеше, сондықтан олардың
құрылымдық және функционалдық қасиеттерінде айырмашылық туындайды.
Зәйтүн майы негізінен моноқанықпаған май қышқылдарына, әсіресе олеин
қышқылына бай (70-78\%). Бұл оның тотығуға төзімділігін арттырып,
технологиялық өңдеу кезінде тұрақтылық береді. Олеин қышқылдарының
жоғары мөлшері эмульсиялардың беріктігін жақсартады, май глобулаларының
дисперсиясын теңестіреді және өнімнің текстуралық қасиеттерін жұмсарта
түседі. Сонымен қатар, құрамындағы табиғи полифенолдар мен Е дәрумені
антиоксидант қасиет көрсетіп, майдың бұзылу жылдамдығын төмендетеді
{[}10, 11{]}.

Күнбағыс майында керісінше полиқанықпаған май қышқылдары басым, әсіресе
линол қышқылы 60-70\% мөлшерде кездеседі. Бұл оның биологиялық
құндылығын арттырғанымен, тотығуға табиғи төзімділігін төмендетеді.
Дегенмен, құрамындағы токоферолдар мөлшері зәйтүн майына қарағанда
жоғары, бұл белгілі бір деңгейде тұрақтылықты сақтайды. полиқанықпаған
май мөлшерінің көп болуы эмульсия құрылымын әлсіретуі мүмкін, сондықтан
өнім құрамындағы ақуыз немесе фосфолипидтер сияқты табиғи
эмульгаторлармен бірге қолдану май бөлінуінің алдын алады және шұжық
өндірісінде майдың біркелкі таралуын қамтамасыз етеді {[}12{]}.

Зығыр майы май қышқылдық құрамы бойынша ең ерекше өнімдердің бірі болып
саналады, себебі оның жалпы майының 50-55\%-ын α-линолен (омега-3)
қышқылы құрайды. Бұл майды функционалдық тағамдарда өте құнды етеді,
алайда омега-3 қышқылдарының тотығуға өте сезімтал болуына байланысты
оның тұрақтылығы төмен, сақтау мерзімі қысқа және термиялық өңдеу
кезінде оңай бұзылады {[}13{]}. Зығыр майы нәзік эмульсия құра алады,
бірақ тұрақты болуы үшін қосымша антиоксиданттар немесе
тұрақтандырғыштар қажет.

Сондықтан оны термиялық өңделетін өнімдерге шектеулі мөлшерде ғана
енгізу ұсынылады {[}14{]}. Осылайша, үш майдың да құрылымдық
қасиеттеріндегі айырмашылық тікелей олардың май қышқылдық құрамына
байланысты: зәйтүн майы - жоғары моноқанықпаған майлардың арқасында
тұрақты және технологиялық тұрғыда қолайлы; күнбағыс майы -- линол
қышқылының жоғары мөлшерінің әсерінен тұрақтылығы төмендеу, бірақ
биологиялық құндылығы жоғары; ал зығыр майы -- омега-3 қышқылының ең бай
көзі ретінде бағалы, бірақ өте тез тотығып, термиялық өңдеуге төзімсіз.
Сондықтан бұл майларды тағам өндірісінде қолданғанда олардың химиялық
құрамына сәйкес технологиялық режимдер таңдалады, ал ақуыз-май
жүйелеріне дұрыс үйлестіру өнім сапасын арттыруға мүмкіндік береді
{[}15, 16{]}.

Ақуыз-май эмульсиясының әртүрлі үлгілерінің функционалды, технологиялық
және органолептикалық қасиеттері ескеріле отырып, оның оңтайлы мөлшері
соя сүті 50\%, зәйтүн майы 44\%, соя ақуыз изоляты - 5\% анықталды
(кесте 3). Бұл көрсеткіш эмульсияның тұрақтылығын сақтауға, құрылымын
біркелкілеуге және тағамдық қасиеттерін жақсартуға мүмкіндік береді.
\end{multicols}

\tcap{3-кесте. Ақуыз-май эмульсиясының бақылау үлгісінің рецептурасы}
\begin{longtblr}[
  label = none,
  entry = none,
]{
  cells = {c},
  cells = {font = \small},
  row{1} = {font = \bfseries\small},
  hlines,
  vlines,
}
{Шикізат, кг/100 кг} & № 1 & № 2 & № 3 & № 4 & № 5 \\
Соя сүті, ақуыз 3-4\%                   & 44           & 47           & 50           & 54           & 60           \\
Зәйтүн майы                             & 50           & 47           & 44           & 40           & 34           \\
Соя ақуыз изоляты                       & 5            & 5            & 5            & 5            & 5            \\
Лимон шырыны                            & 1            & 1            & 1            & 1            & 1            \\
Барлығы:                                & 100          & 100          & 100          & 100          & 100          
\end{longtblr}

\begin{multicols}{2}
Ақуыз-май эмульсиясының сапалық көрсеткіштерін жан-жақты бағалау үшін
әртүрлі арақатынаста дайындалған бес бақылау үлгісі зерттелді. Әр үлгіде
компоненттердің мөлшері өзгертіліп, олардың физика-химиялық құрамы,
ақуыз және май мөлшері және көмірсу мөлшері эмульсия тұрақтылығы мен
құрылымдық-технологиялық қасиеттері анықталды. Бұл талдау әр үлгінің
өзара әрекеттесу ерекшеліктерін бағалап, ақуыз-май эмульсиясының қандай
мөлшерде ең үйлесімді қалыптасатынын көрсетеді.

Соя ақуызының суды байланыстыру қабілеті, май бөлшектерін тұрақтандыруы
және эмульсия дисперстілігіне әсері, сондай-ақ зәйтүн майының май
фазасын біркелкі таратуы мен қоспаның нәзіктілігіне ықпалы әр үлгінің
физика-химиялық нәтижелеріне тікелей әсер етті (кесте 4). Бес түрлі
үлгіні салыстыра отырып, ақуыз-май эмульсиясының құрылымы, тұрақтылығы
және функционалдық қасиеттері бойынша оңтайлы құрамды таңдауға мүмкіндік
туды. Бұл қоспаны әрі қарай технологиялық процестерге енгізу үшін негіз
болып табылады.
\end{multicols}

\tcap{4-кесте. Ақуыз-май эмульсиясының физика-химиялық құрамы}
\begin{longtblr}[
  label = none,
  entry = none,
]{
  cells = {c},
  cells = {font = \small},
  row{1} = {font=\bfseries\small},
  hlines,
  vlines,
}
Көрсеткіштер                                                                        & № 1     & № 2     & № 3     & № 4     & № 5     \\
Ақуыздың массалық үлесі, \%                                                         & 20      & 22      & 24      & 26      & 28      \\
Майдың массалық үлесі, \%                                                           & 25,5    & 24,75   & 24      & 23,25   & 22,5    \\
Көмірсудың массалық үлесі, \%                                                       & 15      & 16,5    & 18      & 19,5    & 21      \\
{Энергетикалық құндылығы:\\\labelitemi\hspace{\dimexpr\labelsep+0.5\tabcolsep}Ккал} & 662,75  & 639,93  & 617,10  & 594,28  & 571,45  \\
\labelitemi\hspace{\dimexpr\labelsep+0.5\tabcolsep}кДж                              & 2777,27 & 2678.73 & 2583,18 & 2487,64 & 2392,09 
\end{longtblr}

\begin{multicols}{2}
Зерттеу барысында ақуыз-май эмульсиясының бес түрлі үлгісі салыстырылып,
олардың тағамдық құрамы мен энергиялық құндылығы анықталды. Әр нұсқада
ақуыз, май және көмірсу мөлшері мақсатты түрде өзгертіліп, композицияның
оптималды пропорциясын табу көзделді. Алынған мәліметтер бойынша үлгілер
арасындағы айырмашылықтар айқын байқалды: ақуыз үлесі 20-28 \%
аралығында, май мөлшері 22,5-25,5 \%, ал көмірсулар 15-21 \% аралығында
өзгерді. Осы өзгерістер қоспаның энергетикалық құндылығына да әсер етті
ккал көрсеткіші 571,45-662,75 аралығында болды.

Берілген мәліметтерді талдай отырып, үшінші үлгінің құрамы теңгерімді
екені байқалды. Бұл үлгіде ақуыз (24 \%) және май (24 \%) мөлшерінің тең
болуы эмульсияның тұрақтылығына, құрылымдық біркелкілікке және
технологиялық жағынан қолайлы консистенцияға ықпал етуі мүмкін. Сонымен
қатар көмірсу мөлшерінің 18 \% деңгейінде болуы жалпы жүйенің энергия
балансын сақтауына әсер етеді. Осы көрсеткіштер үшінші үлгіні ақуыз-май
эмульсиясының оңтайлы нұсқасы ретінде қарастыруға мүмкіндік береді.
\end{multicols}

{\bfseries 1 - сурет. Ақуыз-май эмульсиясының технологиялық сұлбасы}

Жүргізілген зерттеулер негізінде соя дәнінен ақуыз-май эмульсиясын алу
технологиясы жасалды (1 сурет). Технологияның артықшылықтары, соя сүті
және соя ақуыз изоляты эмульсияның тұрақтылығын қамтамасыз етеді, зәйтүн
майымен араластыру нәтижесінде тұрақты, біркелкі құрылымды қоспа алуға
мүмкіндік береді. Бұл процесс, өндірістік жабдықтарды қолданумен
жүргізілетіндіктен, өндірістің қауіпсіздігі мен өнімнің сапасы артады.

Сонымен қатар, алынған эмульсия тағамдық өнімдерге (мысалы, шұжық және
ақуыз-май қоспалары) қосылғанда олардың құрылымы тұрақты, дәмі толық,
текстурасы жұмсақ болып, сақтау мерзімі ұзартылады, шығындарды азайтуға
мүмкіндік береді {[}17{]}. Ақуыз-май эмульсиясының функционалдық
көрсеткіштері 5-ші кестеде және 2 суретте ақуыз-май эмульсиясы
үлгілерінің салыстырмалы органолептикалық көрсеткіштерінің диаграммасы
көрсетілген.

\tcap{5-кесте. Ақуыз-май эмульсиясының функционалдық көрсеткіштері}
\begin{longtblr}[
  label = none,
  entry = none,
]{
  cells = {c},
  cells = {font = \small},
  cell{1}{1} = {font=\bfseries\small},
  cell{1}{2} = {font=\bfseries\small},
  hlines,
  vlines,
}
Көрсеткіштер                           & Бақылау үлгісі          & №\textbf{3 үлгі (оңтайлы)} \\
Су байланыстырғыш қабілеті (СБҚ), \%   & 65 ± 2,1                & 78 ± 1,8                   \\
Май сіңіру қабілеті (МСҚ), \%          & 55 ± 1,5                & 72 ± 1,2                   \\
рН                                     & 5,4 ± 0,1               & 5,6 ± 0,05                 \\
Шығым, \%                              & 92 ± 1,0                & 96 ± 0,8                   \\
Жалпы бактерия, КОЕ/г                  & 5x10\textsuperscript{2} & 2x10\textsuperscript{2}    \\
Сақтау мерзімі, 4 °C температурда, күн & 10 ± 2                  & 17 ± 1,5                   
\end{longtblr}

{\bfseries 2 - сурет. Ақуыз-май эмульсиясы үлгілерінің салыстырмалы
органолептикалық көрсеткіштерінің диаграммасы}

\begin{multicols}{2}
2-ші суретте көрсетілген диаграммада бес түрлі ақуыз-май эмульсиясының
органолептикалық көрсеткіштері салыстырылып берілді. Бағалау нәтижелері
бойынша №3-ші үлгі сыртқы түрі, түсі, иісі, дәмі және консистенциясы
жағынан жоғары балл жинап, басқа үлгілермен салыстырғанда неғұрлым жақсы
сенсорлық көрсеткіш көрсетті. Бұл оның ақуыз бен майдың біркелкі
эмульгирленуі тиімді жүргенін, құрылымының тұрақты қалыптасқанын және
дәмдік қасиеттерінің үйлесімді екенін дәлелдейді.

Ақуыз бен майдың таралуы оңтайлы болуы өнім консистенциясының біртекті
әрі тұрақты қалыптасуына әсер ететін негізгі факторлардың бірі болып
табылады. Осы тұрғыдан алғанда, №3-ші үлгінің құрылымдық тұрақтылығы мен
органолептикалық қабылдануы басқа үлгілерге қарағанда айқын басымдыққа
ие болды.

Алынған нәтижелерді салыстырмалы бағалау үлгілер арасындағы
айырмашылықтардың бар екенін көрсетті. №3-ші үлгінің pH мәні басқа
үлгілермен салыстырғанда тұрақты және оңтайлы диапазонда орналасқаны
анықталды. Бұл оның ақуыздық құрылымының жақсы гидратацияланғанын, май
фазасының біркелкі бөлініп, қоспаның эмульсиялық тұрақтылығы жоғары
болғанын көрсетеді.

{\bfseries Қорытынды.} Зерттеу нәтижесінде соя сүті мен зәйтүн майынан
алынған ақуыз-май эмульсиясының оңтайлы құрамы және оны алудың тиімді
технологиясы анықталды. Бес

нұсқаның салыстырмалы талдауы көрсеткендей, ақуыз, май және
көмірсулардың теңестірілген мөлшері ең жоғары құрылымдық-механикалық
қасиеттерге ие болды және органолептикалық көрсеткіштерді~қамтамасыз
етті. Бұл технология өндірісті жеңілдетіп, өнімнің біркелкілігі мен
қауіпсіздігін арттырады.~Алынған ақуыз-май эмульсияны шұжық өндірісінде
қолдану тағамның құрылымын, дәмін жақсартады және сақтау мерзімін
ұзартады. Сонымен қатар, технология шикізатты тиімді пайдалану
арқылы~шығынды азайтуға~мүмкіндік береді. Жұмыстың зерттеу нәтижелері
ақуыз-май эмульсиясын тағам өндірісіне енгізудің~ғылыми-технологиялық
негізін~қалады және оның функционалдық қасиеттерінің жоғарылығын
дәлелдейді.
\end{multicols}

\begin{center}
{\bfseries Әдебиеттер}
\end{center}

\begin{refs}
1. Қазақстан Республикасы Стратегиялық жоспарлау және реформалар
агенттігі. Ұлттық статистика бюросы.2025 жылғы 1 қазанға Қазақстан
Республикасындағы қолда бар дәнді және бұршақ дақылдары.
\href{https://stat.gov.kz/search/index.php?q=\%D0\%B4\%D3\%99\%D0\%BD\%D0\%B4\%D1\%96+&s}{https://stat.gov.kz} Қаралған күні: 07.09.2025.

2. Chen S., Tang C-H. Effect of varying β-conglycinin / glycinin ratios
on the emulsifying and interfacial properties of soybean
proteins.//\href{https://www.researchgate.net/journal/Modern-Food-Science-and-Technology-1673-9078?_tp=eyJjb250ZXh0Ijp7ImZpcnN0UGFnZSI6InB1YmxpY2F0aW9uIiwicGFnZSI6InB1YmxpY2F0aW9uIn19}{Modern
Food Science and Technology}.-2016.-Vol.32(9).- P.62-68. DOI
\href{https://doi.org/10.13982/j.mfst.1673-9078.2016.9.010}{10.13982/j.mfst.1673-9078.2016.9.010}.

3. \href{https://scijournals.onlinelibrary.wiley.com/authored-by/Liu/Conghui}{Conghui
Liu},~\href{https://scijournals.onlinelibrary.wiley.com/authored-by/Wang/Xin}{Xin
Wang} The physicochemical properties and stability of flaxseed oil
emulsions: effects of emulsification methods and the ratio of soybean
protein isolate to soy
lecithin//\href{https://www.scopus.com/pages/publications/85106255803?origin=document-preview-flyout}{Journal
of the Science of Food and Agriculture} -
2021. -Vol.101(15).-P.6407-6416. DOI 10.1002/jsfa.11311.

4. Shao Jiaqi,Yang Jing,Wu Ping, Jin, Weiping, Xiao Junxia. Effect of
soy protein isolates obtained by weak-base synchronized membrane
separation on nanoemulsion stability: A comparison with alkaline
extraction/isoelectric precipitation and salt
extraction//\href{https://www.scopus.com/pages/publications/105016785177?origin=document-preview-flyout}{Food
Chemistry}Article.- 2025.-Vol.495: 146471.
DOI~10.1016/j.foodchem.2025.146471.

5. Bishunu Karki, Buddhi P. Lamsal, David Grewei, Anthony L. Functional
properties of soy protein isolates produced from ultrasonicated defatted
soy
flakes//\href{https://www.researchgate.net/journal/Journal-of-the-American-Oil-Chemists-Society-1558-9331?_tp=eyJjb250ZXh0Ijp7ImZpcnN0UGFnZSI6InB1YmxpY2F0aW9uIiwicGFnZSI6InB1YmxpY2F0aW9uIn19}{Journal
of the American Oil Chemists'{}
Society}-2009.-Vol.86(10).-P.1021-1028. DOI
\href{https://doi.org/10.1007/s11746-009-1433-0?urlappend=?utm_source=researchgate.net&utm_medium=article}{10.1007/s11746-009-1433-0}.

6. Wenjie Xia,Wing Kie Siu, Leonar M.C. Linear and non-linear rheology
of heat-set soy protein gels:Effects of selective proteolysis of
β-conglycinin and glycinin// Food
Hydrocollioids.-2021.-Vol.120(6):106962. DOI
10.1016/j.foodhyd.2021.106962.

7. Joyce Boye, Fatemeh Zare, Alison Pletch. Pulse proteins: Processing,
characterization, functional properties and applications in food and
feed//Food Research international.-2010.-Vol.43(2).-P.414-431. DOI
10.1016/j.foodres.2009.09.003.

8. Dilek Dulgar Altiner,
\href{https://www.researchgate.net/profile/Seyma_Hallac?_tp=eyJjb250ZXh0Ijp7ImZpcnN0UGFnZSI6InB1YmxpY2F0aW9uIiwicGFnZSI6InB1YmxpY2F0aW9uIn19}{Şeyma
Hallaç}. The effect of soy flour and carob flour addition on the
physicochemical, quality,and sensory properties of pasta formulations//
\href{https://www.researchgate.net/journal/International-Journal-of-Agriculture-Environment-and-Food-Sciences-2602-246X?_tp=eyJjb250ZXh0Ijp7ImZpcnN0UGFnZSI6InB1YmxpY2F0aW9uIiwicGFnZSI6InB1YmxpY2F0aW9uIn19}{International
Journal of Agriculture Environment and Food
Sciences}.-2020.-Vol.4(4).-P.406-417. DOI
\href{https://doi.org/10.31015/jaefs.2020.4.3}{10.31015/jaefs.2020.4.3}.

9. Peanut and Soy Protein-Based Emulsion Gels Loaded with Curcumin as a
New Fat Substitute in Sausages: A Comparative Study.//Gels.- 2025.- Vol.
11(1): 62 DOI 10.3390/gels11010062.

10. Tong Liu, Xinyu Zhen,Nan Zheng,Yaomei Ma, Yu Zhang, Hongyu Lei,
Jiaxin Liu, Ruining Zhang, Jun Zhao. Functional properties and
application potential of black soybean meal insoluble

dietary fiber remodeled by ultrasound combined with ultrahigh
pressure.//LWT. -2025.-
\href{https://www.sciencedirect.com/journal/lwt/vol/221/suppl/C}{Vol.
221}: 117594. DOI 10.1016/j.lwt.2025.117594.

11. \href{https://www.researchgate.net/profile/Ioana-Stanciu-3?_tp=eyJjb250ZXh0Ijp7ImZpcnN0UGFnZSI6InB1YmxpY2F0aW9uIiwicGFnZSI6InB1YmxpY2F0aW9uIn19}{Ioana
Stanciu} Technological Process of Obtaining Sunflower Oil.//Current
Approaches in Engineering Research and Technology.-Vol.9.- P.88-111. DOI
10.9734/bpi/caert/v9/2616.

12. \href{https://www.researchgate.net/profile/Ji-Yao-Zhang?_tp=eyJjb250ZXh0Ijp7ImZpcnN0UGFnZSI6InB1YmxpY2F0aW9uIiwicGFnZSI6InB1YmxpY2F0aW9uIn19}{Ji
Yao Zhang},
\href{https://www.researchgate.net/profile/Ellen-Kim-7?_tp=eyJjb250ZXh0Ijp7ImZpcnN0UGFnZSI6InB1YmxpY2F0aW9uIiwicGFnZSI6InB1YmxpY2F0aW9uIn19}{Ellen
Kim},
\href{https://www.researchgate.net/profile/Peter-Lawrence-15?_tp=eyJjb250ZXh0Ijp7ImZpcnN0UGFnZSI6InB1YmxpY2F0aW9uIiwicGFnZSI6InB1YmxpY2F0aW9uIn19}{Peter
Lawrence}, High Oleic Sunflower Oil Alters the Proportions of n‐6 and
n-3 Polyunsaturated Fatty Acids in Young and Aging C57BL/6 Heart
Tissue//The FASEB Journal.-2015.- Vol.29(S1). DOI
\href{https://doi.org/10.1096/fasebj.29.1_supplement.598.11?urlappend=?utm_source=researchgate.net&utm_medium=article}{10.1096/fasebj.29.1\_supplement.598.11}.

13. Ankit Goyal, Vivek Sharma, Neelam Upadhayay, Sandeep Gill, Manvesh
Sihag. Flaxseed and flaxseed oil: An ancient medicine \& modern
functional food//J Food Sci Technol.- 2014.-Vol.51(9).-P.1633-1653. DOI
\href{https://doi.org/10.1007/s13197-013-1247-9}{10.1007/s13197-013-1247-9}.

14. \href{https://www.researchgate.net/scientific-contributions/Selina-Khan-2122547983?_tp=eyJjb250ZXh0Ijp7ImZpcnN0UGFnZSI6InB1YmxpY2F0aW9uIiwicGFnZSI6InB1YmxpY2F0aW9uIn19}{Selina
Khan},
\href{https://www.researchgate.net/profile/Sharmin-Lisa?_tp=eyJjb250ZXh0Ijp7ImZpcnN0UGFnZSI6InB1YmxpY2F0aW9uIiwicGFnZSI6InB1YmxpY2F0aW9uIn19}{SharminAkter
Lisa},
\href{https://www.researchgate.net/profile/Monira-Obaid?_tp=eyJjb250ZXh0Ijp7ImZpcnN0UGFnZSI6InB1YmxpY2F0aW9uIiwicGFnZSI6InB1YmxpY2F0aW9uIn19}{MoniraObaid}
Tocopherol content of vegetable oils/ fats and their oxidative
deterioration during storage.
//\href{https://www.researchgate.net/journal/World-Journal-of-Pharmacy-and-Pharmaceutical-Sciences-2278-4357?_tp=eyJjb250ZXh0Ijp7ImZpcnN0UGFnZSI6InB1YmxpY2F0aW9uIiwicGFnZSI6InB1YmxpY2F0aW9uIn19}{World
Journal of Pharmacy and Pharmaceutical Sciences.-2015.-}
.-Vol.4(4).-P.1537-1548.

15. Ting Lu,Yan Shen, Zi-Xuan Wu, Hong-Kai, Ao Li,Yong-Fu Wang.
Improving the oxidative stability of flaxseed oil with composite
antioxidants comprising gallic acid alkyl ester with appropriate chain
length//LWT.- 2021.-Vol.138(10):110763. DOI 10.1016/j.lwt.2020.110763.

16. \href{https://pubmed.ncbi.nlm.nih.gov/?term=}{Oleg
Shadyro},\href{https://pubmed.ncbi.nlm.nih.gov/?term=}{Anna
Sosnovskaya},\href{https://pubmed.ncbi.nlm.nih.gov/?term=}{Irina
Edimecheva}. Effect of biologically active substances on oxidative
stability of flaxseed oil//J Food Sci Technol.-
2020. -Vol.57(1).-P.243--252. DOI
\href{https://doi.org/10.1007/s13197-019-04054-4}{10.1007/s13197-019-04054-4}.

17. Ашакаева Р.У., Асенова Б.К. Ет өнімдері үшін байытылған ақуызды--
майлы эмульсиясын жасау.// Наука и образование в современном мире:
вызовы ХХІ века: матер.6-й междунар. науч.-практ. конф. -- Нур-Султан,
2020. -С.50- 52.
\end{refs}

\begin{center}
{\bfseries References}
\end{center}

\begin{refs}
1. Қазақстан Республикасы Стратегиялық жоспарлау және реформалар
агенттігі. Ұлттық статистика бюросы.2025 жылғы 1 қазанға Қазақстан
Республикасындағы қолда бар дәнді және бұршақ дақылдары.
\href{https://stat.gov.kz/search/index.php?q=\%D0\%B4\%D3\%99\%D0\%BD\%D0\%B4\%D1\%96+&s}{https://stat.gov.kz} Қаралған күні: 07.09.2025.

2. Chen S., Tang C-H. Effect of varying β-conglycinin / glycinin ratios
on the emulsifying and interfacial properties of soybean
proteins.//\href{https://www.researchgate.net/journal/Modern-Food-Science-and-Technology-1673-9078?_tp=eyJjb250ZXh0Ijp7ImZpcnN0UGFnZSI6InB1YmxpY2F0aW9uIiwicGFnZSI6InB1YmxpY2F0aW9uIn19}{Modern
Food Science and Technology}.-2016.-Vol.32(9).- P.62-68. DOI
\href{https://doi.org/10.13982/j.mfst.1673-9078.2016.9.010}{10.13982/j.mfst.1673-9078.2016.9.010}.

3. \href{https://scijournals.onlinelibrary.wiley.com/authored-by/Liu/Conghui}{Conghui
Liu},~\href{https://scijournals.onlinelibrary.wiley.com/authored-by/Wang/Xin}{Xin
Wang} The physicochemical properties and stability of flaxseed oil
emulsions: effects of emulsification methods and the ratio of soybean
protein isolate to soy
lecithin//\href{https://www.scopus.com/pages/publications/85106255803?origin=document-preview-flyout}{Journal
of the Science of Food and Agriculture} -
2021. -Vol.101(15).-P.6407-6416. DOI 10.1002/jsfa.11311.

4. Shao Jiaqi,Yang Jing,Wu Ping, Jin, Weiping, Xiao Junxia. Effect of
soy protein isolates obtained by weak-base synchronized membrane
separation on nanoemulsion stability: A comparison with alkaline
extraction/isoelectric precipitation and salt
extraction//\href{https://www.scopus.com/pages/publications/105016785177?origin=document-preview-flyout}{Food
Chemistry}Article.- 2025.-Vol.495: 146471.
DOI~10.1016/j.foodchem.2025.146471.

5. Bishunu Karki, Buddhi P. Lamsal, David Grewei, Anthony L. Functional
properties of soy protein isolates produced from ultrasonicated defatted
soy flakes//\href{https://www.researchgate.net/journal/Journal-of-the-American-Oil-Chemists-Society-1558-9331?_tp=eyJjb250ZXh0Ijp7ImZpcnN0UGFnZSI6InB1YmxpY2F0aW9uIiwicGFnZSI6InB1YmxpY2F0aW9uIn19}{Journal
of the American Oil Chemists'{}
Society}-2009.-Vol.86(10).-P.1021-1028. DOI
\href{https://doi.org/10.1007/s11746-009-1433-0?urlappend=?utm_source=researchgate.net&utm_medium=article}{10.1007/s11746-009-1433-0}.

6. Wenjie Xia,Wing Kie Siu, Leonar M.C. Linear and non-linear rheology
of heat-set soy protein gels:Effects of selective proteolysis of
β-conglycinin and glycinin// Food
Hydrocollioids.-2021.-Vol.120(6):106962. DOI
10.1016/j.foodhyd.2021.106962.

7. Joyce Boye, Fatemeh Zare, Alison Pletch. Pulse proteins: Processing,
characterization, functional properties and applications in food and
feed//Food Research international.-2010.-Vol.43(2).-P.414-431. DOI
10.1016/j.foodres.2009.09.003.

8. Dilek Dulgar Altiner,
\href{https://www.researchgate.net/profile/Seyma_Hallac?_tp=eyJjb250ZXh0Ijp7ImZpcnN0UGFnZSI6InB1YmxpY2F0aW9uIiwicGFnZSI6InB1YmxpY2F0aW9uIn19}{Şeyma
Hallaç}. The effect of soy flour and carob flour addition on the
physicochemical, quality,and sensory properties of pasta formulations//
\href{https://www.researchgate.net/journal/International-Journal-of-Agriculture-Environment-and-Food-Sciences-2602-246X?_tp=eyJjb250ZXh0Ijp7ImZpcnN0UGFnZSI6InB1YmxpY2F0aW9uIiwicGFnZSI6InB1YmxpY2F0aW9uIn19}{International
Journal of Agriculture Environment and Food
Sciences}.-2020.-Vol.4(4).-P.406-417. DOI
\href{https://doi.org/10.31015/jaefs.2020.4.3}{10.31015/jaefs.2020.4.3}.

9. Peanut and Soy Protein-Based Emulsion Gels Loaded with Curcumin as a
New Fat Substitute in Sausages: A Comparative Study.//Gels.- 2025.- Vol.
11(1): 62 DOI 10.3390/gels11010062.

10. Tong Liu, Xinyu Zhen,Nan Zheng,Yaomei Ma, Yu Zhang, Hongyu Lei,
Jiaxin Liu, Ruining Zhang, Jun Zhao. Functional properties and
application potential of black soybean meal insoluble
dietary fiber remodeled by ultrasound combined with ultrahigh
pressure.//LWT. -2025.-
\href{https://www.sciencedirect.com/journal/lwt/vol/221/suppl/C}{Vol.
221}: 117594. DOI 10.1016/j.lwt.2025.117594.

11. \href{https://www.researchgate.net/profile/Ioana-Stanciu-3?_tp=eyJjb250ZXh0Ijp7ImZpcnN0UGFnZSI6InB1YmxpY2F0aW9uIiwicGFnZSI6InB1YmxpY2F0aW9uIn19}{Ioana
Stanciu} Technological Process of Obtaining Sunflower Oil.//Current
Approaches in Engineering Research and Technology.-Vol.9.- P.88-111. DOI
10.9734/bpi/caert/v9/2616.

12. \href{https://www.researchgate.net/profile/Ji-Yao-Zhang?_tp=eyJjb250ZXh0Ijp7ImZpcnN0UGFnZSI6InB1YmxpY2F0aW9uIiwicGFnZSI6InB1YmxpY2F0aW9uIn19}{Ji
Yao Zhang},
\href{https://www.researchgate.net/profile/Ellen-Kim-7?_tp=eyJjb250ZXh0Ijp7ImZpcnN0UGFnZSI6InB1YmxpY2F0aW9uIiwicGFnZSI6InB1YmxpY2F0aW9uIn19}{Ellen
Kim},
\href{https://www.researchgate.net/profile/Peter-Lawrence-15?_tp=eyJjb250ZXh0Ijp7ImZpcnN0UGFnZSI6InB1YmxpY2F0aW9uIiwicGFnZSI6InB1YmxpY2F0aW9uIn19}{Peter
Lawrence}, High Oleic Sunflower Oil Alters the Proportions of n‐6 and
n‐3 Polyunsaturated Fatty Acids in Young and Aging C57BL/6 Heart
Tissue//The FASEB Journal.-2015.- Vol.29(S1). DOI
\href{https://doi.org/10.1096/fasebj.29.1_supplement.598.11?urlappend=?utm_source=researchgate.net&utm_medium=article}{10.1096/fasebj.29.1\_supplement.598.11}.

13. Ankit Goyal, Vivek Sharma, Neelam Upadhayay, Sandeep Gill, Manvesh
Sihag. Flaxseed and flaxseed oil: An ancient medicine \& modern
functional food//J Food Sci Technol.- 2014.-Vol.51(9).-P.1633-1653. DOI
\href{https://doi.org/10.1007/s13197-013-1247-9}{10.1007/s13197-013-1247-9}.

14. \href{https://www.researchgate.net/scientific-contributions/Selina-Khan-2122547983?_tp=eyJjb250ZXh0Ijp7ImZpcnN0UGFnZSI6InB1YmxpY2F0aW9uIiwicGFnZSI6InB1YmxpY2F0aW9uIn19}{Selina
Khan},
\href{https://www.researchgate.net/profile/Sharmin-Lisa?_tp=eyJjb250ZXh0Ijp7ImZpcnN0UGFnZSI6InB1YmxpY2F0aW9uIiwicGFnZSI6InB1YmxpY2F0aW9uIn19}{SharminAkter
Lisa},
\href{https://www.researchgate.net/profile/Monira-Obaid?_tp=eyJjb250ZXh0Ijp7ImZpcnN0UGFnZSI6InB1YmxpY2F0aW9uIiwicGFnZSI6InB1YmxpY2F0aW9uIn19}{MoniraObaid}
Tocopherol content of vegetable oils/ fats and their oxidative
deterioration during storage.
//\href{https://www.researchgate.net/journal/World-Journal-of-Pharmacy-and-Pharmaceutical-Sciences-2278-4357?_tp=eyJjb250ZXh0Ijp7ImZpcnN0UGFnZSI6InB1YmxpY2F0aW9uIiwicGFnZSI6InB1YmxpY2F0aW9uIn19}{World
Journal of Pharmacy and Pharmaceutical Sciences.-2015.-}
.-Vol.4(4).-P.1537-1548.

15. Ting Lu,Yan Shen, Zi-Xuan Wu, Hong-Kai, Ao Li,Yong-Fu Wang.
Improving the oxidative stability of flaxseed oil with composite
antioxidants comprising gallic acid alkyl ester with appropriate chain
length//LWT.- 2021.-Vol.138(10):110763. DOI 10.1016/j.lwt.2020.110763.

16. \href{https://pubmed.ncbi.nlm.nih.gov/?term=}{Oleg
Shadyro},\href{https://pubmed.ncbi.nlm.nih.gov/?term=}{Anna
Sosnovskaya},\href{https://pubmed.ncbi.nlm.nih.gov/?term=}{Irina
Edimecheva}. Effect of biologically active substances on oxidative
stability of flaxseed oil//J Food Sci Technol.-
2020. -Vol.57(1).-P.243--252. DOI
\href{https://doi.org/10.1007/s13197-019-04054-4}{10.1007/s13197-019-04054-4}.

17. Ashakaeva R.U., Asenova B.K. Et өnіmderі үshіn bajytylғan aқuyzdy -
majly jemul' sijasyn zhasau.// Nauka i obrazovanie v
sovremennom mire: vyzovy HHІ veka: mater.6-j mezhdunar. nauch.-prakt.
konf. - Nur-Sultan, 2020. -S.50- 52. {[}in Kazakh{]}
\end{refs}

\begin{info}
\hspace{1em}\emph{{\bfseries Авторлар туралы мәліметтер}}

Байтукенова Ш.Б. - т.ғ.к., қауымдастырылған профессор, «С. Сейфуллин
атындағы Қазақ агротехникалық зерттеу университеті» КеАҚ, Астана,
Қазақстан, e-mail:
baytukenova75@mail.ru;

Бақытбек М. - магистрант «С. Сейфуллин атындағы Қазақ агротехникалық
зерттеу университеті» КеАҚ, Астана, Қазақстан, e-mail:
mila9231115@gmail.com;

Байтукенова С.Б. - т.ғ.к., қауымдастырылған профессор, «Қ.Құлажанов
атындағы Қазақ технология және бизнес университеті» АҚ, Астана қ.,
Казахстан, e-mail: saule7272 @mail.ru,

Костанова А.Т. - докторант НАО «Казахский агротехнический
исследовательский университет им. С.Сейфуллина», Астана, Казахстан,
e-mail: anel\_kostanova@mail.ru.

\hspace{1em}\emph{{\bfseries Information about the authors}}

Baitukenova Sh.B. - Candidate of technical sciences, associate
professor, "S.Seifullin Kazakh Agrotechnical Research University",
Astana, Kazakhstan, e-mail: baytukenova75@mail.ru;

Bakytbek M. - Master' s student of the «The Department of
Food Technology and Processing Products», «S.Seifullin Kazakh
Agrotechnical Research University», Astana, Kazakhstan, e-mail:
mila9231115@gmail.com;

Baitukenova S.B. - Candidate of technical sciences, associate professor,
JSC "Kazakh University of technology and business named after K.
Kulazhanov", Kazakhstan, Astana, e-mail: saule7272 @mail.ru;

Kostanova A.Т. - PhD student of the «The Department of Food Technology
and Processing Products», «S.Seifullin Kazakh Agrotechnical Research
University», Astana, Kazakhstan, e-mail: anel\_kostanova@mail.ru.
\end{info}
